\documentclass[11pt,a4paper]{article}
% preamble.tex — estilo común de todos los apuntes Froth.
% Cada apunte hace % preamble.tex — estilo común de todos los apuntes Froth.
% Cada apunte hace % preamble.tex — estilo común de todos los apuntes Froth.
% Cada apunte hace \input{preamble} justo después de \documentclass.
\usepackage[margin=2.4cm]{geometry}
\usepackage{xcolor}
\definecolor{litblue}{RGB}{31,79,121}
\definecolor{litgreen}{RGB}{27,94,32}
\definecolor{litorange}{RGB}{230,81,0}
\usepackage{parskip}
\usepackage{titlesec}
\titleformat{\section}{\Large\bfseries\color{litblue}}{\thesection}{0.6em}{}
\titleformat{\subsection}{\large\bfseries\color{litblue}}{\thesubsection}{0.6em}{}
\usepackage{tcolorbox}
\usepackage{fancyhdr}
\pagestyle{fancy}
\fancyhf{}
\fancyhead[L]{\small\color{litblue}Froth — Apuntes de ML}
\fancyhead[R]{\small\thepage}
\renewcommand{\headrulewidth}{0.4pt}
\usepackage[colorlinks=true,linkcolor=litblue,urlcolor=litblue]{hyperref}

% Cabecera de cada apunte: \cabecera{numero}{titulo}
\newcommand{\cabecera}[2]{%
  \begin{tcolorbox}[colback=litblue,coltext=white,colframe=litblue,arc=2mm]
    {\Large\bfseries Apunte #1 — #2}\\[3pt]
    {\small Proyecto Froth · aprender Machine Learning construyendo}
  \end{tcolorbox}\vspace{0.8em}}

% Cajas temáticas
\newtcolorbox{concepto}[1]{colback=blue!4,colframe=litblue,title={Concepto: #1},arc=1.5mm}
\newtcolorbox{practica}{colback=green!5,colframe=litgreen,title={Pruébalo tú (teclado en mano)},arc=1.5mm}
\newtcolorbox{ojo}{colback=orange!8,colframe=litorange,title={Ojo / error típico},arc=1.5mm}
\newtcolorbox{resumen}{colback=gray!8,colframe=gray!60,title={En una frase},arc=1.5mm}

\newcommand{\codigo}[1]{\texttt{#1}}
 justo después de \documentclass.
\usepackage[margin=2.4cm]{geometry}
\usepackage{xcolor}
\definecolor{litblue}{RGB}{31,79,121}
\definecolor{litgreen}{RGB}{27,94,32}
\definecolor{litorange}{RGB}{230,81,0}
\usepackage{parskip}
\usepackage{titlesec}
\titleformat{\section}{\Large\bfseries\color{litblue}}{\thesection}{0.6em}{}
\titleformat{\subsection}{\large\bfseries\color{litblue}}{\thesubsection}{0.6em}{}
\usepackage{tcolorbox}
\usepackage{fancyhdr}
\pagestyle{fancy}
\fancyhf{}
\fancyhead[L]{\small\color{litblue}Froth — Apuntes de ML}
\fancyhead[R]{\small\thepage}
\renewcommand{\headrulewidth}{0.4pt}
\usepackage[colorlinks=true,linkcolor=litblue,urlcolor=litblue]{hyperref}

% Cabecera de cada apunte: \cabecera{numero}{titulo}
\newcommand{\cabecera}[2]{%
  \begin{tcolorbox}[colback=litblue,coltext=white,colframe=litblue,arc=2mm]
    {\Large\bfseries Apunte #1 — #2}\\[3pt]
    {\small Proyecto Froth · aprender Machine Learning construyendo}
  \end{tcolorbox}\vspace{0.8em}}

% Cajas temáticas
\newtcolorbox{concepto}[1]{colback=blue!4,colframe=litblue,title={Concepto: #1},arc=1.5mm}
\newtcolorbox{practica}{colback=green!5,colframe=litgreen,title={Pruébalo tú (teclado en mano)},arc=1.5mm}
\newtcolorbox{ojo}{colback=orange!8,colframe=litorange,title={Ojo / error típico},arc=1.5mm}
\newtcolorbox{resumen}{colback=gray!8,colframe=gray!60,title={En una frase},arc=1.5mm}

\newcommand{\codigo}[1]{\texttt{#1}}
 justo después de \documentclass.
\usepackage[margin=2.4cm]{geometry}
\usepackage{xcolor}
\definecolor{litblue}{RGB}{31,79,121}
\definecolor{litgreen}{RGB}{27,94,32}
\definecolor{litorange}{RGB}{230,81,0}
\usepackage{parskip}
\usepackage{titlesec}
\titleformat{\section}{\Large\bfseries\color{litblue}}{\thesection}{0.6em}{}
\titleformat{\subsection}{\large\bfseries\color{litblue}}{\thesubsection}{0.6em}{}
\usepackage{tcolorbox}
\usepackage{fancyhdr}
\pagestyle{fancy}
\fancyhf{}
\fancyhead[L]{\small\color{litblue}Froth — Apuntes de ML}
\fancyhead[R]{\small\thepage}
\renewcommand{\headrulewidth}{0.4pt}
\usepackage[colorlinks=true,linkcolor=litblue,urlcolor=litblue]{hyperref}

% Cabecera de cada apunte: \cabecera{numero}{titulo}
\newcommand{\cabecera}[2]{%
  \begin{tcolorbox}[colback=litblue,coltext=white,colframe=litblue,arc=2mm]
    {\Large\bfseries Apunte #1 — #2}\\[3pt]
    {\small Proyecto Froth · aprender Machine Learning construyendo}
  \end{tcolorbox}\vspace{0.8em}}

% Cajas temáticas
\newtcolorbox{concepto}[1]{colback=blue!4,colframe=litblue,title={Concepto: #1},arc=1.5mm}
\newtcolorbox{practica}{colback=green!5,colframe=litgreen,title={Pruébalo tú (teclado en mano)},arc=1.5mm}
\newtcolorbox{ojo}{colback=orange!8,colframe=litorange,title={Ojo / error típico},arc=1.5mm}
\newtcolorbox{resumen}{colback=gray!8,colframe=gray!60,title={En una frase},arc=1.5mm}

\newcommand{\codigo}[1]{\texttt{#1}}

\begin{document}
\cabecera{21}{El mapa mental de temas (y el hallazgo de los silos)}

\section{Subir un nivel: de papers a temas}

La red del Apunte 20 tiene 120 nodos — buena para explorar, mala para \emph{presentar}.
El \textbf{mapa mental de temas} agrega la información un nivel arriba: cada
\textbf{subtema entero se vuelve UN nodo} (tamaño $=$ nº de papers) y las aristas resumen
la relación entre subtemas. Es la vista ejecutiva: 3--4 burbujas que cuentan todo el campo.

\begin{concepto}{agregación: la misma red, menos resolución}
Este truco —agrupar nodos y resumir sus relaciones— se llama \emph{agregación} y es
universal: de personas $\rightarrow$ departamentos, de calles $\rightarrow$ ciudades, de
papers $\rightarrow$ subtemas. Se pierde detalle, se gana legibilidad. Por eso mantenemos
\emph{ambas} redes: la de papers (explorar) y la de temas (entender/presentar).
\end{concepto}

\section{Las aristas llevan DOS números (y ahí está la gracia)}

Entre cada par de subtemas medimos dos cosas distintas:
\begin{enumerate}
  \item \textbf{Similitud} $=$ promedio del coseno entre los papers de uno y otro (en 768D).
        ``¿Qué tan relacionado es lo que \emph{estudian}?''
  \item \textbf{Citas cruzadas} $=$ cuántas veces un paper de un subtema \emph{cita} a uno
        del otro. ``¿Qué tanto se \emph{leen} entre ellos?''
\end{enumerate}

\section{El hallazgo: silos (tu gap, cuantificado)}

Con tus datos salió esto:

\begin{center}
\begin{tabular}{lcc}
\textbf{Relación} & \textbf{Similitud} & \textbf{Citas cruzadas} \\
\hline
Geología $\leftrightarrow$ Flotación & 0.688 & \textbf{1} \\
Geología $\leftrightarrow$ Economía/reciclaje & 0.662 & 3 \\
Flotación $\leftrightarrow$ Economía/reciclaje & 0.683 & 1 \\
\end{tabular}
\end{center}

\begin{ojo}
Lee la primera fila con ojos de tesista: geología y flotación \textbf{estudian cosas muy
relacionadas} (0.69 es alto) pero \textbf{casi no se citan} — ¡1 cita entre 90 papers! De
las 116 citas internas del corpus, solo 5 cruzan fronteras de subtema: las comunidades son
\textbf{silos} — conectadas por significado, separadas por costumbre.

Tu tesis (``flotación de micas de Li \emph{y su impacto en} la recuperación de by-products
del granito'') vive exactamente en ese silencio entre silos. Esto es un \emph{gap}
cuantificado y citable en tu introducción: la Fase 5 formalizará esta detección.
\end{ojo}

\section{Dónde vive en el código}

\codigo{froth/graph.py} $\rightarrow$ \codigo{build\_topic\_graph()}: nodos con
label/nº papers/citas/centroide, aristas con similitud y citas cruzadas. Se imprime al
correr \codigo{python -m froth.graph}; la vista previa está en
\codigo{3\_Resultados/topic\_mindmap.png}.

\begin{resumen}
Agregar la red de papers a nivel de temas da la vista ejecutiva (3--4 burbujas). Cada arista
lleva similitud (qué estudian) y citas (qué se leen) — y en tu corpus divergen: subtemas muy
afines que no se citan $=$ silos $=$ tu gap, ya con números.
\end{resumen}

\end{document}
